\chapter{Plan de verificación experimental detallado}\label{cap:ensayos}

Este capítulo convierte el plan E1--E6 del capítulo~\ref{cap:riesgos} en procedimientos
ejecutables: qué banco se arma, qué se mide, con qué incertidumbre y con qué criterio se
acepta cada requisito. Agrega cinco ensayos (E7--E11) para cubrir los requisitos que el
plan original no tocaba, y separa E5 en vibración (E5a) y choque (E5b). Todos los valores son \textbf{predicciones o estimaciones}; ningún
ensayo se ha hecho todavía.

\begin{resultado}[Resumen del plan]
\begin{itemize}[leftmargin=*]
  \item \textbf{Banco vertical} (E2, E3): chorro abierto de Ø\,\SI{0.8}{m} o más
        ($D_j\ge2D$), \SIrange{6}{8}{m/s}. Mide $\CR$ con
        $U\approx\pm\num{0.055}$ (95\,\%), rpm, conicidad y par.
  \item \textbf{Auto} (E2b): cubre lo que el ventilador no alcanza, el arranque a
        \SIrange{10}{14}{m/s}, con tramos cortos para no pasar de las rpm de diseño.
  \item \textbf{Caídas desde drone} (E4): de 80 a \SI{100}{m} de caída por ensayo según el
        modelo nominal (los \SIrange{30}{60}{m} del capítulo~\ref{cap:riesgos} no alcanzan);
        10 caídas dan $\Vr$ con $U\approx\pm\SI{0.33}{m/s}$. Si la dispersión entre caídas
        es de \SI{0.4}{m/s} o menos, 3 caídas alcanzan para demostrar $\Vr\le\SI{7.5}{m/s}$.
  \item \textbf{Ambientales} (E5, E7--E9) con recursos de universidad, y una matriz que
        lleva cada requisito del anexo a su método y criterio.
\end{itemize}
\end{resultado}

%=====================================================================
\section{Magnitudes a medir}\label{sec:ensayos-magnitudes}
El diseño descansa sobre cinco números: el coeficiente $\CR$, la velocidad estabilizada
$\Vr$, las rpm, la conicidad $\beta_0$ y el par de arranque $Q_0$. La
tabla~\ref{tab:ensayos-magnitudes} fija para cada uno el valor esperado, la incertidumbre
que hace falta para decidir y el ensayo que lo mide.

\begin{table}[H]
\centering\small
\caption{Magnitudes clave, valor esperado e incertidumbre objetivo ($k\approx2$).}
\label{tab:ensayos-magnitudes}
\begin{tabularx}{\textwidth}{@{}llcYc@{}}
\toprule
\textbf{Magnitud} & \textbf{Esperado} & \textbf{$U$ objetivo} & \textbf{Para qué se necesita} & \textbf{Ensayo}\\\midrule
$\CR$ del rotor & 1,2 & $\pm5$\,\% & $\Vr$ por análisis; $\pm5$\,\% en $\CR$ es $\pm2{,}5$\,\% en $\Vr$ & E3\\
$\Vr$ con drogue & \SI{6.4}{m/s} & \SI{\pm0.3}{m/s} & Cumplimiento directo de \Req{C4} & E4\\
rpm de equilibrio & \num{1150} & $\pm2$\,\% & Validar la ecuación~\eqref{eq:omega} y el $C_l$ de la placa & E2\\
Conicidad $\beta_0$ & \ang{4.4} & $\pm\ang{0.3}$ & Validar la ecuación~\eqref{eq:conicidad}; margen al tope de \ang{10} & E2\\
Par de arranque $Q_0$ & \SI{0.026}{N.m} & $\pm20$\,\% & Elegir el paso; riesgo n.º~1 del AMFE & E2, E2b\\
Tiempo de arranque & \SIrange{3}{6}{s} & \SI{\pm0.5}{s} & Altura perdida antes de frenar & E2b, E4\\
\bottomrule
\end{tabularx}
\end{table}

%=====================================================================
\section{Banco vertical de ventilador (E2 y E3)}\label{sec:ensayos-banco}

\subsection{Tamaño del chorro y bloqueo}
El chorro debe tener al menos el doble del diámetro del rotor. A $\Vr/\vh=1{,}83$ el rotor
está en el estado de estela turbulenta (figura~\ref{fig:estados}): el caudal neto a
través del disco es casi nulo y el rotor se comporta como un cuerpo romo de
$\CR\approx1{,}2$. Por eso no sirve la corrección clásica de Glauert para hélices~\cite{ensayos-barlow1999},
que supone un tubo de corriente bien definido a través del disco, hipótesis que falla en
este estado. El modelo adecuado es el de bloqueo de cuerpos
romos.

\paragraph{Cámara cerrada.} Si se usa el túnel de la universidad, Maskell da la presión
dinámica corregida~\cite{ensayos-maskell1963}
\begin{equation}
  \frac{q_c}{q} = 1 + \theta\,\CR\,\frac{A_R}{C}, \qquad \theta\approx2{,}5,
  \label{eq:ensayos-maskell}
\end{equation}
con $C$ el área de la sección de ensayo. La teoría se dedujo para placas macizas; aplicada
a un rotor es una aproximación. Con secciones habituales la corrección es grande
(tabla~\ref{tab:ensayos-opciones}): en un túnel de \qtyproduct{1 x 1}{m} la presión
dinámica efectiva sería un 38\,\% mayor que la leída, y una incertidumbre del 20\,\% sobre esa
corrección ya invalida el $\pm5$\,\% buscado.

\paragraph{Chorro abierto.} La frontera libre permite que el chorro se ensanche; la
corrección es menor y de signo opuesto a la de cámara cerrada, pero mezcla expansión y
desvío del chorro, bloqueo de la boquilla y gradiente de presión
estática~\cite{ensayos-mercker1996}. Como ningún modelo cubre bien un rotor en estela
turbulenta, \textbf{la corrección se mide}: antes del rotor se ensaya un disco macizo
delgado de Ø\,\SI{400}{m m} en la misma posición. Su coeficiente en aire libre es
$C_D\approx1{,}17$~\cite{hoerner1965}, con una dispersión de algunos por ciento entre
fuentes; el factor $F_b=1{,}17/C_{D,\text{medido}}$ se aplica
luego al rotor, que tiene casi el mismo coeficiente y por lo tanto casi el mismo bloqueo.

\begin{table}[H]
\centering\small
\caption{Opciones de banco para el rotor de Ø\,\SI{400}{mm}. Potencia eléctrica con
rendimiento global de 0,35 (estimación).}
\label{tab:ensayos-opciones}
\begin{tabular}{@{}lcccc@{}}
\toprule
\textbf{Instalación} & $\boldsymbol{A_R/A_j}$ \textbf{o} $\boldsymbol{A_R/C}$ & \textbf{Corrección} & \textbf{$P$ a \SI{7.8}{m/s}} & \textbf{$P$ a \SI{13.4}{m/s}}\\\midrule
Chorro abierto Ø\,\SI{0.6}{m}   & 0,44 & no admisible  & \SI{0.23}{kW} & \SI{1.2}{kW}\\
\textbf{Chorro abierto Ø\,\SI{0.8}{m}} & \textbf{0,25} & \textbf{medida con disco} & \textbf{\SI{0.42}{kW}} & \textbf{\SI{2.1}{kW}}\\
Chorro abierto Ø\,\SI{1.0}{m}   & 0,16 & medida con disco & \SI{0.65}{kW} & \SI{3.3}{kW}\\
Túnel cerrado \qtyproduct{1.0 x 1.0}{m} & 0,126 & $q_c/q=1{,}38$ & --- & ---\\
Túnel cerrado \qtyproduct{1.5 x 1.5}{m} & 0,056 & $q_c/q=1{,}17$ & --- & ---\\
Túnel cerrado \qtyproduct{2.0 x 2.0}{m} & 0,031 & $q_c/q=1{,}09$ & --- & ---\\
\bottomrule
\end{tabular}
\end{table}

\begin{resultado}[Banco elegido]
Chorro abierto vertical de Ø\,\SIrange{0.8}{1.0}{m} hecho con un ventilador industrial de
tambor (\SIrange{0.5}{0.7}{kW}) y un variador. Cubre el punto de operación
(\SIrange{6}{8}{m/s}). La velocidad de arranque (\SI{13.4}{m/s}) pediría más de
\SI{2}{kW}: ese caso pasa al ensayo en auto (sección~\ref{sec:ensayos-auto}).
\end{resultado}

\subsection{Calidad del flujo}
El chorro debe ser uniforme y sin remolino; si no, las rpm medidas quedan sesgadas. Un
remolino de ángulo $\psi$ suma $V\tan\psi$ a la velocidad tangencial de la pala: a
\SI{6.4}{m/s} y $\psi=\ang{2}$ son \SI{0.22}{m/s} frente a $\Omega r=\SI{18}{m/s}$ en
$0{,}75R$, es decir un sesgo de 1,2\,\% en las rpm. El acondicionamiento sigue las
reglas de Mehta y Bradshaw~\cite{ensayos-mehta1979}: panal de celdas con largo de 6 a 8
veces su diámetro (elimina el remolino) y dos mallas finas (bajan la turbulencia).
\begin{itemize}
  \item \textbf{Uniformidad}: mapa con Pitot sin rotor, en 13 puntos (centro y dos anillos
        de 6 puntos en $r=\SIlist{100;180}{mm}$) en el plano del rotor. Criterio: $q$ dentro
        de $\pm2$\,\% de la media en Ø\,\SI{440}{mm}. El mismo mapa da el factor
        $K=\bar q_\text{disco}/q_\text{ref}$ entre la presión media en el disco y la de la
        sonda de referencia.
  \item \textbf{Remolino}: hilos de lana en una grilla; criterio $|\psi|<\ang{2}$.
  \item \textbf{Posición}: el núcleo uniforme de un chorro redondo se angosta con la
        distancia (del orden de 4 a 5 diámetros de largo); el rotor va a
        $0{,}5$--$1\,D_j$ de la boquilla, dentro del núcleo y fuera de la zona de bloqueo
        de la boquilla.
\end{itemize}

\subsection{Instrumentación}
La figura~\ref{fig:ensayos-banco} muestra el banco. La maqueta del cuerpo (Ø\,95) va
\textbf{aguas arriba} del rotor, como en vuelo, con soporte propio que no carga la
balanza. El rotor cuelga de un vástago que pasa por el eje hueco, en la posición de la
cuerda del drogue, y lleva el empuje a la celda de carga.

\begin{table}[H]
\centering\small
\caption{Instrumentación del banco. Las incertidumbres son objetivos tras la calibración.}
\label{tab:ensayos-instr}
\begin{tabularx}{\textwidth}{@{}lYcc@{}}
\toprule
\textbf{Magnitud} & \textbf{Instrumento y calibración} & \textbf{Rango} & \textbf{$u$ objetivo}\\\midrule
Presión dinámica $q$ & Pitot-estática Ø\,\SI{3}{mm} + SDP810-125Pa (cero \SI{0.08}{Pa}; ganancia 3\,\% sin calibrar~\cite{ensayos-sdp810}), calibrado con manómetro inclinado & \SI{\pm125}{Pa} & 0,5\,\%\\
Velocidad (control) & Anemómetro de hélice; solo verificación cruzada & \SIrange{0.5}{20}{m/s} & 2--3\,\%\\
Empuje $T$ & Celda de viga de \SI{2}{kg} + HX711 a \SI{80}{Hz}; masas patrón de \SIrange{50}{1000}{g} & \SI{20}{N} & 0,3\,\%\\
rpm & Hall de a bordo (1 imán) y tacómetro óptico (4 pulsos por vuelta) & \num{0}--\num{3000} & 0,5\,\%\\
Par $Q$ & Aceleración (ec.~\ref{eq:ensayos-Qacel}) o freno de Foucault con celda de \SI{100}{g} & \SI{0.05}{N.m} & 5\,\%\\
Conicidad $\beta_0$ & Cámara lateral, \SI{240}{fps}, exposición 1/2000\,s & \ang{0}--\ang{12} & \ang{0.1}\\
$p$, $T$ ambiente & Barómetro y termómetro de referencia para $\rho$ & --- & 0,3\,\% en $\rho$\\
\bottomrule
\end{tabularx}
\end{table}

\begin{figure}[H]
\centering
\begin{tikzpicture}[x=0.74mm, y=0.74mm,
    lab/.style={font=\scriptsize, align=left},
    labr/.style={font=\scriptsize, align=right},
    ld/.style={ref}]
  % bastidor y piso
  \draw[cgris, line width=1.2pt] (18,0) -- (18,120) -- (132,120) -- (132,0);
  \draw[black, line width=0.8pt] (-2,0) -- (152,0);
  % ventilador
  \draw[fijo] (35,1) rectangle (115,12);
  \node[font=\scriptsize] at (75,6.5) {ventilador de tambor + variador};
  % cámara de estabilización, panal y mallas
  \draw[cfijo, line width=0.7pt] (35,12) -- (35,40); \draw[cfijo, line width=0.7pt] (115,12) -- (115,40);
  \fill[pattern=crosshatch, pattern color=cgris!70] (35,17) rectangle (115,23);
  \draw[cgris] (35,17) rectangle (115,23);
  \draw[cgris, densely dotted, line width=0.9pt] (35,29) -- (115,29);
  \draw[cgris, densely dotted, line width=0.9pt] (35,34) -- (115,34);
  % contracción hasta la boquilla (D_j = 60)
  \draw[cfijo, line width=0.7pt] (35,40) .. controls (35,50) and (45,46) .. (45,55);
  \draw[cfijo, line width=0.7pt] (115,40) .. controls (115,50) and (105,46) .. (105,55);
  % capas de mezcla
  \draw[cgris, dashed] (45,55) -- (53,108);
  \draw[cgris, dashed] (105,55) -- (97,108);
  \draw[cgris!50, dashed] (45,55) -- (38,108);
  \draw[cgris!50, dashed] (105,55) -- (112,108);
  \foreach \x in {52,60,90,98} \draw[flecha, cfijo!55] (\x,57) -- (\x,68);
  % maqueta del cuerpo (aguas arriba) y su soporte
  \draw[fijo] (71.5,60) rectangle (78.5,79);
  \draw[cgris, line width=0.9pt] (71.5,64) -- (18,64);
  % rotor con conicidad
  \draw[rot] (72.5,80) rectangle (77.5,84.5);
  \draw[crot, line width=1.5pt] (60,84.2) -- (72.5,82.4);
  \draw[crot, line width=1.5pt] (77.5,82.4) -- (90,84.2);
  \fill[black] (77.5,81) circle (0.7);
  % vástago y celda de carga
  \draw[cgris, line width=1.1pt] (75,84.5) -- (75,106);
  \draw[black, fill=white, line width=0.7pt] (68,106) rectangle (82,112);
  \draw[black, line width=0.9pt] (75,112) -- (75,120);
  % Pitot de referencia (punta hacia abajo, contra el flujo)
  \draw[black, line width=0.9pt] (100,56) -- (100,61) -- (112,61);
  \draw[black, fill=white] (112,58) rectangle (121,64);
  % tacómetro óptico
  \draw[black, fill=white] (99,91) rectangle (107,96);
  \draw[cfreno, densely dashed] (99,92) -- (86,84);
  % cámara lateral
  \draw[black, fill=white] (140,80) rectangle (150,87);
  \draw[black, fill=white] (140,81.5) -- (136,80) -- (136,87) -- (140,85.5);
  \draw[cgris, densely dotted] (136,83.5) -- (91,83.5);
  % adquisición
  \node[caja, font=\scriptsize, minimum width=22mm] (daq) at (150,30) {Adquisición\\ HX711, SDP810,\\ hall, PC};
  \draw[cgris] (82,109) -- (146,109) -- (146,40);
  \draw[cgris] (121,61) -- (150,61) -- (150,40);
  % cota distancia boquilla-rotor
  \draw[cota] (127,55) -- (127,83.5);
  \draw[ld] (106,55) -- (129,55);
  \node[font=\scriptsize, rotate=90, anchor=south] at (126,74) {$0{,}5$--$1\,D_j$};
  % etiquetas izquierda
  \node[labr, anchor=east] (l1) at (14,20) {panal\\ $L/d=6$--$8$};
  \draw[ld] (14,20) -- (35,20);
  \node[labr, anchor=east] at (14,31.5) {2 mallas};
  \draw[ld] (14,31.5) -- (35,31.5);
  \node[labr, anchor=east] at (14,47) {contracción};
  \draw[ld] (14,47) -- (39,47);
  \node[labr, anchor=east] at (14,72) {maqueta del\\ cuerpo Ø\,95,\\ no métrica};
  \draw[ld] (14,72) -- (71.5,72);
  \node[labr, anchor=east] at (14,98) {chorro abierto\\ $D_j\ge2D$};
  \draw[ld] (14,98) -- (40,98);
  \node[labr, anchor=east] at (14,112) {celda de carga\\ \SI{2}{kg}};
  \draw[ld] (14,112) -- (68,109);
  % etiquetas derecha
  \node[lab, anchor=west] at (136,103) {vástago por el\\ eje hueco};
  \draw[ld] (136,103) -- (75,103);
  \node[lab, anchor=west] at (108,93.5) {tacómetro\\ óptico};
  \node[lab, anchor=west] at (136,74) {cámara lateral\\ \SI{240}{fps}};
  \node[lab, anchor=north west] at (90,81.5) {rotor Ø\,400};
  \node[lab, anchor=south] at (116.5,64.2) {Pitot};
  \node[lab, anchor=west] at (79,79) {hall};
\end{tikzpicture}
\caption{Banco vertical de chorro abierto (esquema, sin escala). Naranja: rotor; azul:
partes fijas; rojo punteado: haz del tacómetro óptico.}
\label{fig:ensayos-banco}
\end{figure}


\paragraph{Velocidad.} A \SI{6}{m/s} la presión dinámica es apenas \SI{22}{Pa}, y el
cero de \SI{0.08}{Pa} pesa un 0,36\,\% de $q$. El rango de \SI{125}{Pa} llega a
\SI{14.3}{m/s}. La sonda de \SI{3}{mm} trabaja a $Re\approx\num{1200}$, donde los efectos
viscosos sobre el Pitot son despreciables. El coeficiente se calcula directamente con $q$,
sin pasar por $\rho$:
\begin{equation}
  \CR = \frac{F_b\,T}{K\,q_\text{ref}\,A_R}, \qquad
  \lambda = \frac{\Omega R}{V},\quad V=\sqrt{2Kq_\text{ref}/\rho}.
  \label{eq:ensayos-CR}
\end{equation}

\paragraph{Par.} Dos métodos independientes:
\begin{enumerate}
  \item \textbf{Aceleración.} Durante el arranque con el chorro a velocidad fija,
  \begin{equation}
    Q_\text{aero}(\Omega,V) = I_p\,\frac{d\Omega}{dt} + Q_f(\Omega),
    \label{eq:ensayos-Qacel}
  \end{equation}
  donde $Q_f$ es el rozamiento, medido con un ensayo de desaceleración libre sin palas. La
  derivada sale de los períodos del sensor hall, suavizados con un filtro de
  Savitzky--Golay. Una aceleración de \SI{10}{rad/s^2} ya equivale a \SI{5.1e-3}{N.m}.
  La inercia $I_p$ se mide con un péndulo trifilar,
  \begin{equation}
    I = \frac{m\,g\,r^2\,\tau^2}{4\pi^2 L} - I_\text{plato},
    \label{eq:ensayos-trifilar}
  \end{equation}
  con hilos de $L=\SI{0.8}{m}$ a $r=\SI{0.10}{m}$: el período es $\tau\approx\SI{1.5}{s}$ y
  se toma el promedio de 20 oscilaciones.
  \item \textbf{Freno de Foucault.} Un disco de aluminio de \SI{2}{mm} en el cubo y un imán
  de neodimio sobre un brazo que apoya en una celda de \SI{100}{g}. Al acercar el imán
  aumenta el par de freno y se recorre la curva $Q(\Omega)$ completa a $V$ fija. La
  pendiente $\partial Q/\partial\Omega$ en el equilibrio es el dato que más falta en el
  modelo de arranque (sección~\ref{sec:ensayos-drone}).
\end{enumerate}

\paragraph{Conicidad.} La punta sube $\Delta z=(R-e)\sin\beta_0\approx\SI{12}{mm}$ sobre la
bisagra, de donde
\begin{equation}
  \beta_0 = \arcsin\frac{\Delta z}{R-e}.
  \label{eq:ensayos-beta}
\end{equation}
Con \SI{0.25}{mm} por píxel y $\pm2$ píxeles de error de borde,
$u(\beta_0)\approx\ang{0.1}$. La cámara va a \SI{2}{m} o más y con zoom, para que la
perspectiva no incline el plano de las puntas.

\subsection{Matriz y procedimiento}
\begin{enumerate}
  \item Sin rotor: mapa de $q$ (13 puntos), factor $K$, remolino; disco de referencia para
        $F_b$ a \SIlist{6;8}{m/s}.
  \item Tara: cubo y vástago sin palas, a cada velocidad.
  \item Para cada paso ($\ang{-2}$, $\ang{-3}$, $\ang{-4}$): velocidades
        \SIlist{5;6;6.5;7;8}{m/s}; en cada una se esperan rpm estables
        ($|\dot\Omega|<\SI{1}{rad/s^2}$ durante \SI{5}{s}) y se promedian \SI{10}{s} de
        $T$, $q$ y rpm. Cada punto se repite 5 veces desmontando y volviendo a montar.
  \item Arranques: 10 por paso, con el rotor frenado a mano y liberado con el chorro
        ya establecido. Se registra la curva $\Omega(t)$.
  \item Video lateral a \SIlist{6;8}{m/s} para $\beta_0$; foto con flash para el
        seguimiento de puntas (diferencia de altura entre palas $<\SI{2}{mm}$).
\end{enumerate}
Criterio de E3: $\CR\ge1{,}0$ con su cota inferior al 95\,\%. El punto de diseño se acepta
si $\CR\ge1{,}1$; con $\CR$ entre 1,0 y 1,1 el diseño cumple \Req{C4}, pero con menos
margen (tabla~\ref{tab:rotor}).

%=====================================================================
\section{Alternativa en auto (E2b)}\label{sec:ensayos-auto}
El auto da gratis el rango de \SIrange{10}{14}{m/s} (\SIrange{36}{50}{km/h}), que el
ventilador no alcanza; por eso se usa para el arranque, no para medir $\CR$. A velocidad
constante el rotor no se queda en las rpm de diseño: tiende a la misma relación de
velocidad de punta $\lambda_d\approx3{,}76$, es decir a unas \num{2400}~rpm a
\SI{13.4}{m/s} (punta a \SI{50}{m/s}, fuerza centrífuga de \SI{57}{N} por pala, 4,4 veces
la de diseño). En vuelo eso no ocurre porque el CanSat frena mientras el rotor acelera
(figura~\ref{fig:ensayos-caida}: máximo de unas \num{1220}~rpm). Cada tramo en auto mide
entonces solo el arranque, de 0 a \num{1150}~rpm, y al llegar a ese valor el auto baja a
\SI{25}{km/h} (unas \num{1250}~rpm de equilibrio). El rotor va
con el eje horizontal, alineado con la marcha, en un mástil a \SI{1}{m} o más sobre el
techo (figura~\ref{fig:ensayos-auto}).

\begin{itemize}
  \item \textbf{Velocidad}: se mide la del aire con el Pitot del mástil, no la del
        velocímetro ni la del GPS. Cada condición se corre en los dos sentidos del mismo
        tramo y se promedian los tiempos de arranque, lo que cancela casi todo el viento
        longitudinal.
  \item \textbf{Viento cruzado}: produce flujo oblicuo. Una veleta con potenciómetro, o
        un hilo filmado, mide el ángulo. Se descartan tramos con más de \ang{5}, que a
        \SI{13.4}{m/s} equivalen a \SI{1.2}{m/s} de viento cruzado. Por eso se ensaya
        con calma, temprano por la mañana.
  \item \textbf{Gravedad}: con el eje horizontal el peso de cada pala
        (\SI{0.073}{N}) actúa en el plano de giro. Es el 0,6\,\% de la fuerza centrífuga
        en régimen, pero en el arranque las palas cuelgan de forma asimétrica. Se espera que
        eso alargue algo el arranque (no está cuantificado), así que los tiempos en auto
        se toman como una cota pesimista.
  \item \textbf{Flujo del vehículo}: el techo acelera el aire. El Pitot va a la misma
        altura que el rotor y \SI{0.5}{m} al costado, y mide ese efecto en lugar de
        corregirlo.
\end{itemize}

\begin{nota}[Seguridad en el ensayo en auto]
Solo en un predio cerrado (pista, aeródromo o estacionamiento vacío), nunca en la vía
pública; máximo \SI{50}{km/h}. Mástil sujeto a las barras de techo con cintas con
trinquete; malla de protección alrededor del disco. Nadie en el plano de giro: una pala
suelta de \SI{7.4}{g} lleva \SI{0.8}{J} a \num{1150}~rpm y \SI{3.5}{J} a \num{2400}~rpm
(velocidad de su centro de masa, a \SI{123}{mm} del eje). Antes del primer tramo, cada
portapala se prueba con una carga estática de \SI{60}{N} (la centrífuga a \num{2400}~rpm).
Si el hall pasa de \num{1500}~rpm, el auto frena. El conductor solo conduce; un
segundo operador maneja la adquisición, y antes de cada tramo se revisan pasadores y
portapalas.
\end{nota}

\begin{figure}[H]
\centering
\begin{subfigure}[b]{0.52\textwidth}\centering
\begin{tikzpicture}[x=0.85mm, y=0.85mm]
  \draw[fijo] (0,8) -- (0,18) -- (12,20) -- (22,30) -- (52,30) -- (60,20) -- (74,18) -- (74,8) -- cycle;
  \draw[black, fill=white] (14,8) circle (5); \draw[black, fill=white] (60,8) circle (5);
  \draw[black, line width=0.8pt] (-4,3) -- (82,3);
  \draw[cgris, line width=1.2pt] (30,30.8) -- (52,30.8);
  % mástil, maqueta del cuerpo aguas arriba, cubo y disco
  \draw[cgris, line width=1.3pt] (48,30.8) -- (48,56) -- (40,56);
  \draw[fijo] (24,54.5) rectangle (35,57.5);
  \draw[rot] (35,54) rectangle (40,58);
  \draw[crot, line width=1.5pt] (40.5,44) -- (40.5,68);
  % Pitot
  \draw[black, line width=0.9pt] (48,40) -- (12,40);
  \fill[black] (12,40) circle (0.8);
  % cotas y flechas
  \draw[cota] (66,30.8) -- (66,44);
  \draw[ref] (41,44) -- (68,44);
  \node[font=\scriptsize, anchor=west] at (67,37.4) {$\ge\SI{1}{m}$};
  \draw[flecha, cfijo] (-2,64) -- (20,64) node[midway, above, font=\scriptsize] {viento relativo};
  \draw[flecha] (40,-1.5) -- (20,-1.5) node[midway, below, font=\scriptsize] {marcha};
  \node[font=\scriptsize, anchor=east, align=right] at (11,40) {Pitot y\\ veleta};
  \node[font=\scriptsize, anchor=west] at (42,67) {rotor};
\end{tikzpicture}
\caption{E2b: mástil sobre el auto.}
\end{subfigure}\hfill
\begin{subfigure}[b]{0.44\textwidth}\centering
\begin{tikzpicture}[x=0.85mm, y=0.85mm]
  \fill[pattern=north east lines, pattern color=cgris] (0,70) rectangle (60,74);
  \draw[cgris, line width=0.8pt] (0,70) -- (60,70);
  \draw[black, line width=0.8pt] (20,70) -- (20,68.2); \draw[black, line width=0.7pt] (20,67) circle (1.2);
  % cuerda tensa y CanSat abajo
  \draw[ccuerda, line width=1.2pt] (20,65.8) -- (20,30);
  \draw[fijo] (16,8) rectangle (24,30);
  \node[marca] at (20,19) {a};
  % posición inicial
  \draw[fijo, dashed, fill=none] (34,48) rectangle (42,70);
  \draw[ccuerda, densely dashed] (20,65.8) .. controls (28,55) and (36,72) .. (38,70);
  \draw[flecha, cgris] (47,66) -- (47,32);
  \node[font=\scriptsize, anchor=west, align=left] at (48,49) {caída\\ libre\\ \SI{0.61}{m}};
  \draw[cota] (8,66) -- (8,30);
  \node[font=\scriptsize, rotate=90, anchor=south] at (7,48) {cuerda \SI{61}{cm}};
  \node[font=\scriptsize, anchor=west, align=left] at (26,14) {acelerómetro\\ $\pm\SI{200}{g}$, \SI{1}{kHz}};
  \draw[black, line width=0.8pt] (0,2) -- (60,2);
\end{tikzpicture}
\caption{E5b: choque por caída con cuerda~\cite{ensayos-cansatenv}.}
\end{subfigure}
\caption{Esquemas de los ensayos en auto y de choque (sin escala).}
\label{fig:ensayos-auto}
\end{figure}

%=====================================================================
\section{Caídas desde drone (E4)}\label{sec:ensayos-drone}

\subsection{Altura necesaria}
Hacen falta de 80 a \SI{100}{m} de caída por ensayo si la pendiente del par en el
equilibrio es la estimada con la teoría de elemento de pala (caso nominal), y más de
\SI{120}{m} si es mucho menor (caso conservador). Los \SIrange{30}{60}{m} previstos en
el capítulo~\ref{cap:riesgos} no alcanzan para llegar a régimen. Para estimarlo se
integra un modelo de orden reducido con tres estados (altura, velocidad y rpm):
\begin{align}
  m\,\dot V &= W - (k_d+k_c)V^2 - T, &
  T &= \tfrac12\rho A_R\Bigl[C_s V^2 + (\CR - C_s)\Bigl(\frac{\Omega R}{\lambda_d}\Bigr)^2\Bigr],
  \label{eq:ensayos-modeloV}\\
  I_p\,\dot\Omega &= Q - Q_f, &
  Q &= Q_0\!\left(\frac{V}{13{,}4}\right)^{\!2}(1-x)(1+a\,x),\quad x=\frac{\Omega R}{\lambda_d V}.
  \label{eq:ensayos-modeloQ}
\end{align}
Aquí $C_s=0{,}27$ es el coeficiente con las palas abiertas y quietas,
$\lambda_d=3{,}76$ la relación de velocidad de punta de diseño y $Q_0=\SI{0.026}{N.m}$ el
par de arranque de la ecuación~\eqref{eq:Q0} con paso \ang{-2};
$Q_f=\SI{3e-4}{N.m}$ es una estimación del rozamiento. El parámetro $a$ fija la pendiente
del par en el equilibrio. En autorrotación el par de perfil
$Q_p=B\tfrac12\rho c\,C_{d0}\Omega^2(R^4-e^4)/4\approx\SI{0.038}{N.m}$ (con $C_{d0}=0{,}06$)
crece con $\Omega^2$, mientras que el par de sustentación crece menos; así resulta
$\partial Q/\partial\Omega\approx-k_Q\,Q_p/\Omega$ con $k_Q$ entre 1 y 2. El caso
\textbf{nominal} usa $k_Q=1{,}5$ ($a=8{,}6$). El caso \textbf{conservador} usa $a=0$, un
par lineal que cae desde $Q_0$, unas 10 veces menos pendiente en el equilibrio.

\begin{figure}[H]
\centering
\begin{tikzpicture}
\begin{axis}[width=0.95\textwidth, height=5.8cm, xmin=0, xmax=150, ymin=0, ymax=15,
  xlabel={Altura caída desde la suelta [m]}, ylabel={Velocidad de descenso $V$ [m/s]},
  unbounded coords=jump,
  legend columns=2, legend style={font=\footnotesize, at={(0.5,-0.30)}, anchor=north, cells={anchor=west}, /tikz/every even column/.append style={column sep=8pt}}]
  \fill[cfreno!8] (axis cs:122,0) rectangle (axis cs:150,15);
  \node[font=\footnotesize, cfreno, anchor=north east, align=right] at (axis cs:149,12.6) {sobre\\ \SI{122}{m}};
  \addplot[crot, very thick] table[x=h, y=VAnom] {analysis/data/ensayos_caida.dat};
  \addlegendentry{A: secuencia real, nominal}
  \addplot[crot, thick, dashed] table[x=h, y=VAcons] {analysis/data/ensayos_caida.dat};
  \addlegendentry{A: secuencia real, conservador}
  \addplot[cfijo, very thick] table[x=h, y=VBnom] {analysis/data/ensayos_caida.dat};
  \addlegendentry{B: rotor libre desde la suelta, nominal}
  \addplot[cfijo, thick, dashed] table[x=h, y=VBcons] {analysis/data/ensayos_caida.dat};
  \addlegendentry{B: rotor libre, conservador}
  \addplot[ccuerda, thick] table[x=h, y=VCnom] {analysis/data/ensayos_caida.dat};
  \addlegendentry{C: rotor libre sin drogue, nominal}
  \draw[black!60, densely dotted] (axis cs:0,6.4) -- (axis cs:150,6.4);
  \node[font=\scriptsize, anchor=north west] at (axis cs:90,6.3) {$\Vr=\SI{6.4}{m/s}$};
  \draw[black!60, densely dotted] (axis cs:0,13.35) -- (axis cs:150,13.35);
  \node[font=\scriptsize, anchor=north west] at (axis cs:60,13.25) {$V_1=\SI{13.4}{m/s}$};
  \draw[cgris, line width=0.6pt, {Latex[length=1.6mm]}-{Latex[length=1.6mm]}] (axis cs:57,2.0) -- (axis cs:89,2.0);
  \node[font=\scriptsize, anchor=south] at (axis cs:73,2.0) {ventana de \SI{5}{s} (caso A)};
\end{axis}
\end{tikzpicture}
\caption{Velocidad en función de la altura caída, según el modelo de orden
reducido~\eqref{eq:ensayos-modeloV}--\eqref{eq:ensayos-modeloQ}. En el caso A el rotor se
libera a los \SI{2.5}{s} ($V\approx\SI{12.7}{m/s}$). Zona roja: por encima del límite de
la categoría abierta~\cite{ensayos-anac2025}.}
\label{fig:ensayos-caida}
\end{figure}

\begin{table}[H]
\centering\small
\caption{Presupuesto de altura por caída. $h_{\text{2\,\%}}$: altura caída hasta quedar a
$\pm2$\,\% de la velocidad final; se suman \SI{5}{s} de ventana de medición y \SI{10}{m}
de margen.}
\label{tab:ensayos-altura}
\begin{tabular}{@{}llccc@{}}
\toprule
& & \multicolumn{2}{c}{\textbf{Modelo nominal}} & \textbf{Conservador}\\\cmidrule(lr){3-4}\cmidrule(l){5-5}
\textbf{Caso} & \textbf{Configuración} & $\boldsymbol{h}_{\text{\textbf{2\,\%}}}$ & \textbf{Altura necesaria} & \textbf{Altura necesaria}\\\midrule
A & Secuencia real: drogue y luego rotor & \SI{57}{m} & \SI{99}{m} & \SI{186}{m}\\
B & Rotor libre desde la suelta, con drogue & \SI{39}{m} & \SI{81}{m} & \SI{169}{m}\\
C & Rotor libre, sin drogue & \SI{43}{m} & \SI{89}{m} & \SI{135}{m}\\
\bottomrule
\end{tabular}
\end{table}

\begin{nota}[Límite de altura y permisos]
La categoría abierta de la ANAC permite volar a la vista hasta \SI{122}{m} sobre el
terreno~\cite{ensayos-anac2025}. Antes de la campaña hay que confirmar si soltar una carga
desde el drone exige autorización de categoría específica. La suelta se hace sobre un
campo despejado y vallado, sin personas debajo.
\end{nota}

Consecuencias:
\begin{itemize}
  \item Las caídas de caracterización (repeticiones para $\Vr$ y $\CR$) se hacen en la
        configuración B, la más corta, soltando desde \SIrange{110}{120}{m}.
  \item La secuencia real A (drogue, liberación, arranque) se hace 3 veces desde la
        altura máxima. Además de $\Vr$, verifica la lógica de liberación. No sirve para
        \Req{C3}: a los \SI{2.5}{s} (\SI{21}{m}) la etapa 1 va a \SI{12.7}{m/s}, el 95\,\%
        de $V_1$, todavía acelerando. $V_1$ se mide en caídas aparte con el rotor trabado,
        que llegan al 98\,\% de $V_1$ tras unos \SI{29}{m}.
  \item Si el rotor real se parece al caso conservador, a \SI{120}{m} no llega a
        régimen. Esto se detecta a bordo: la ventana de análisis solo es válida si el
        sensor hall muestra $|\dot\Omega|<\SI{2}{rad/s^2}$. Un arranque tan lento sería
        además un hallazgo de diseño, porque también alarga la caída real, y llevaría a un
        paso más negativo.
  \item \Req{C4} pide un \emph{promedio} en todo el descenso, no la velocidad
        estabilizada. Con liberación a $V_1$ estable, el modelo da una media de la etapa 2
        de \SI{6.51}{m/s} en \SI{560}{m} (apogeo de \SI{700}{m}) en el caso nominal, y de
        \SI{6.86}{m/s} en el conservador; en \SI{100}{m}, de \SI{6.98}{m/s} y
        \SI{8.41}{m/s}. El transitorio pesa poco en vuelo, pero domina en una caída corta,
        por eso en E4 se mide $\Vr$ en la ventana estable y la media de vuelo se obtiene
        por análisis.
\end{itemize}

\subsection{Suspensión y secuencia}
El CanSat cuelga del liberador \textbf{por el vértice del drogue}, con el drogue estirado y
la cuerda tensa; así se espera que se infle en menos de \SI{1}{s} tras la suelta (el
modelo supone \SI{0.5}{s}). En el caso A la
liberación del rotor la ordena el software con un umbral de tiempo, o el comando
\texttt{MEC} (capítulo~\ref{cap:elec}). En el caso B la traba se abre en tierra y las palas
se sujetan con un hilo que corta el mismo liberador. Cada caída registra a bordo el
barómetro a \SI{50}{Hz}, la IMU, el sensor hall y el video de las dos cámaras. En tierra se
anotan la temperatura, la presión y el viento a \SI{2}{m} (anemómetro de mano).

\subsection{Análisis de los datos del barómetro}
La velocidad sale de la presión, no de la altitud ISA:
\begin{equation}
  V = \frac{1}{\rho\,g}\,\frac{dp}{dt}, \qquad \rho = \frac{p}{R_a\,T_\text{loc}},
  \label{eq:ensayos-baro}
\end{equation}
con $V$ positiva hacia abajo (al descender la presión crece) y la temperatura medida en
el lugar. Usar la atmósfera ISA ($\SI{15}{\celsius}$) un día
de \SI{25}{\celsius} sesga la velocidad en un 3,5\,\%. La pendiente se ajusta por
mínimos cuadrados en la ventana estable. Con $N$ muestras de ruido $\sigma_h$ a frecuencia
$f_s$,
\begin{equation}
  \sigma_V = \sigma_h\,f_s\sqrt{\frac{12}{N(N^2-1)}},
  \label{eq:ensayos-sigmaV}
\end{equation}
que con $\sigma_h=\SI{0.3}{m}$, \SI{50}{Hz} y \SI{5}{s} da solo \SI{0.013}{m/s}. El ruido
no es el problema. Las fuentes que dominan son otras:
\begin{itemize}
  \item \textbf{Presión estática de la toma.} La presión dinámica pasa de \SI{109}{Pa} en
        la etapa 1 a \SI{25}{Pa} en la etapa 2. Con un coeficiente de presión de 0,3 a 0,6
        en la toma, eso produce un salto aparente de \SIrange{2}{4}{m} de altura en la
        transición. Por eso no se usa la ventana de transición. La toma va en el costado
        del cuerpo, con varios orificios repartidos en la circunferencia, lejos del rotor y
        de la cara inferior: en descenso esa cara mira al flujo y es un punto de remanso
        ($C_p\approx1$), que daría un salto de unos \SI{7}{m}.
  \item \textbf{Viento vertical.} Una térmica de \SI{0.5}{m/s} cambia $V$ en un 8\,\%.
        Se vuela temprano, sobre terreno uniforme, y se toma la media de muchas caídas.
  \item \textbf{Deriva meteorológica}: \SI{1}{hPa} en \SI{3}{h} equivale a
        \SI{0.015}{m} en \SI{20}{s}, despreciable.
\end{itemize}
La velocidad se informa normalizada al nivel del mar ISA,
$V_\text{ISA}=V\sqrt{\rho/\rho_0}$, para poder compararla con el diseño.

\subsection{Número de repeticiones}\label{sec:ensayos-n}
El número de caídas lo fija la dispersión entre caídas $s$, que depende del viento
vertical y no se conoce de antemano. La semiamplitud del intervalo al 95\,\% de la media es
\begin{equation}
  E = t_{0{,}975;\,n-1}\,\frac{s}{\sqrt n}.
  \label{eq:ensayos-E}
\end{equation}

\begin{figure}[H]
\centering
\begin{subfigure}[b]{0.49\textwidth}\centering
\begin{tikzpicture}
\begin{axis}[width=0.97\textwidth, height=5.4cm, xmin=3, xmax=30, ymin=0, ymax=1.0,
  xlabel={Número de caídas $n$}, ylabel={$E$ al 95\,\% [m/s]},
  legend style={font=\scriptsize, at={(0.97,0.97)}, anchor=north east, cells={anchor=west}}]
  \addplot[cfijo, thick] table[x=n, y=s02] {analysis/data/ensayos_ic.dat}; \addlegendentry{$s=0{,}2$}
  \addplot[ccuerda, thick] table[x=n, y=s03] {analysis/data/ensayos_ic.dat}; \addlegendentry{$s=0{,}3$}
  \addplot[crot, thick] table[x=n, y=s04] {analysis/data/ensayos_ic.dat}; \addlegendentry{$s=0{,}4$}
  \addplot[cfreno, thick] table[x=n, y=s05] {analysis/data/ensayos_ic.dat}; \addlegendentry{$s=0{,}5$}
  \addplot[black!60, dashed, domain=3:30] {0.25}; \addlegendentry{\SI{\pm0.25}{m/s}}
  \addplot[black!60, dotted, thick, domain=3:30] {0.16}; \addlegendentry{$\pm2{,}5$\,\% de $\Vr$}
\end{axis}
\end{tikzpicture}
\caption{Semiamplitud del intervalo de $\Vr$.}
\label{fig:ensayos-ic}
\end{subfigure}\hfill
\begin{subfigure}[b]{0.49\textwidth}\centering
\begin{tikzpicture}
\begin{semilogxaxis}[width=0.97\textwidth, height=5.4cm, xmin=500, xmax=5e5, ymin=0, ymax=120,
  xlabel={Rigidez equivalente $k$ [N/m]}, ylabel={Pico $G$ (fuerza/peso)},
  xticklabel style={font=\scriptsize}]
  \addplot[cfijo, very thick] table[x=k, y=G] {analysis/data/ensayos_choque.dat};
  \draw[cfreno, dashed] (axis cs:500,30) -- (axis cs:5e5,30);
  \draw[cfreno, dashed] (axis cs:3376,0) -- (axis cs:3376,30);
  \node[font=\scriptsize, cfreno, anchor=north west, align=left, fill=white, inner sep=1pt] at (axis cs:4500,26) {30\,G con\\ $k\approx\SI{3.4}{kN/m}$};
  \node[font=\scriptsize, anchor=north west, align=left] at (axis cs:600,115) {cuerda sola de Kevlar:\\ cientos de G (ideal)};
\end{semilogxaxis}
\end{tikzpicture}
\caption{Choque de la caída con cuerda de \SI{0.61}{m}.}
\label{fig:ensayos-G}
\end{subfigure}
\caption{(a) Caídas necesarias según la dispersión $s$ [m/s]. (b) Pico de aceleración
en función de la rigidez de toda la cadena cuerda--soporte--CanSat
(ec.~\ref{eq:ensayos-G}).}
\label{fig:ensayos-icG}
\end{figure}

\begin{resultado}[Plan secuencial de caídas]
Se hacen 3 caídas y se calcula $s$. Para el \textbf{cumplimiento} basta la cota superior
unilateral al 95\,\%, con el límite de \SI{7.5}{m/s} que fija el equipo como margen sobre
los \SI{8}{m/s} del anexo: si $\bar V=\SI{6.4}{m/s}$, con $s=\SI{0.4}{m/s}$ y $n=3$ es
$\bar V+t_{0{,}95;2}\,s/\sqrt3=\SI{7.07}{m/s}<\SI{7.5}{m/s}$. Para la
\textbf{caracterización} se completan las caídas que pida el $s$ medido
(figura~\ref{fig:ensayos-ic}): con $s=\SI{0.3}{m/s}$ hacen falta 9 caídas para
$\pm\SI{0.25}{m/s}$ y 16 para $\pm2{,}5$\,\%; con $s=\SI{0.4}{m/s}$, 13 y 27. Medir
$\CR$ al $\pm5$\,\% por caídas es caro; para eso está el banco.
\end{resultado}

%=====================================================================
\section{Incertidumbre de $\CR$ y de $\Vr$}\label{sec:ensayos-gum}
La incertidumbre se evalúa según la GUM~\cite{ensayos-gum2008}: incertidumbre típica
combinada $u_c^2=\sum(c_i u_i)^2$, grados de libertad efectivos por Welch--Satterthwaite y
factor de cobertura $k=t_{0{,}975;\nu_\text{eff}}$. Los valores de $u_i$ son
\textbf{supuestos de planificación}; tras los ensayos se reemplazan por los medidos.

\begin{table}[H]
\centering\footnotesize
\caption{Presupuesto de incertidumbre de $\CR$ en el banco
(ec.~\ref{eq:ensayos-CR}), a \SI{6.4}{m/s}. Todos los coeficientes de sensibilidad
relativos valen $\pm1$.}
\label{tab:ensayos-gumCR}
\begin{tabular}{@{}llccr@{}}
\toprule
\textbf{Fuente} & \textbf{Evaluación} & \textbf{Distribución} & $\boldsymbol{u_i}$ \textbf{relativa} & $\boldsymbol{\nu_i}$\\\midrule
$T$: calibración de la celda con masas patrón & B & normal & 0,30\,\% & 50\\
$T$: ruido promediado en \SI{10}{s} & A & normal & 0,11\,\% & 30\\
$T$: tara del cubo y del vástago (\SI{5}{mN}) & B & normal & 0,13\,\% & 10\\
$q_\text{ref}$: sensor calibrado contra manómetro & B & normal & 0,50\,\% & 50\\
$q_\text{ref}$: cero (\SI{0.08}{Pa}) & B & rectangular & 0,18\,\% & $\infty$\\
$q_\text{ref}$: fluctuación del chorro & A & normal & 0,30\,\% & 30\\
$K$: no uniformidad del chorro & A & normal & 1,00\,\% & 12\\
$A_R$: radio $\pm\SI{0.5}{mm}$ & B & rectangular & 0,29\,\% & $\infty$\\
$F_b$: bloqueo y $C_D$ del disco ($\pm3$\,\%) & B & rectangular & 1,73\,\% & 8\\
Repetibilidad (5 montajes, $s=1$\,\%) & A & normal & 0,45\,\% & 4\\\midrule
\multicolumn{3}{@{}l}{\textbf{Combinada} $u_c$; $\nu_\text{eff}=19$; $k=2{,}09$} & \textbf{2,19\,\%} & \\
\multicolumn{3}{@{}l}{\textbf{Expandida} $U$ (95\,\%)} & \textbf{4,6\,\%} & \\
\bottomrule
\end{tabular}
\end{table}

\begin{table}[H]
\centering\small
\caption{Presupuesto de incertidumbre de $\Vr$ (con drogue) por caídas desde drone, con
$n=10$ y $s=\SI{0.35}{m/s}$ supuestos.}
\label{tab:ensayos-gumV}
\begin{tabular}{@{}llcr@{}}
\toprule
\textbf{Fuente} & \textbf{Evaluación} & $\boldsymbol{u_i}$ [m/s] & $\boldsymbol{\nu_i}$\\\midrule
Dispersión entre caídas, $s/\sqrt n$ & A & 0,111 & 9\\
Viento vertical medio de la campaña (no se promedia) & B & 0,100 & 8\\
Transitorio residual en la ventana & B & 0,030 & 8\\
Escala barométrica ($T_\text{loc}$ con $\pm\SI{1}{K}$) & B & 0,022 & $\infty$\\
Presión estática de la toma & B & 0,020 & 8\\
Ganancia del sensor de presión ($\pm0{,}5$\,\%, rectangular) & B & 0,019 & $\infty$\\
Normalización a densidad ISA & B & 0,010 & $\infty$\\
Masa ($\pm\SI{1}{g}$) & B & 0,002 & $\infty$\\\midrule
\textbf{Combinada} $u_c$; $\nu_\text{eff}=20$; $k=2{,}08$ & & \textbf{0,157} & \\
\textbf{Expandida} $U$ (95\,\%) & & \textbf{0,33} & \\
\bottomrule
\end{tabular}
\end{table}

\begin{resultado}[Incertidumbre esperada]
\begin{itemize}[leftmargin=*]
  \item Banco: $\CR=1{,}20\pm0{,}055$ (95\,\%). Domina el bloqueo del chorro (1,7\,\%), y
        por eso se mide con el disco de referencia.
  \item Caídas: $\Vr=\SI{6.40 \pm 0.33}{m/s}$ (95\,\%). Dominan la dispersión y el viento
        vertical; más caídas solo reducen la primera.
  \item $\CR$ deducido de caídas \emph{sin} drogue: $u=4{,}6$\,\%, el doble que en el
        banco, porque $\CR\propto V^{-2}$ y además entra el $C_D$ del cuerpo. La
        comparación entre los dos valores de $\CR$ mide la interacción rotor--cuerpo.
        La comparación de $\Vr$ con y sin drogue mide cuánto aporta el drogue en la estela
        del rotor.
\end{itemize}
\end{resultado}

%=====================================================================
\section{Ensayos ambientales y de sistema}\label{sec:ensayos-amb}

\subsection{Vibración de 15\,G (E5a)}
Un vibrador electrodinámico de \SI{200}{N} o más alcanza los 15\,G entre \SIlist{50;500}{Hz}.
A 15\,G de pico la amplitud de desplazamiento es \SI{1.5}{mm} a \SI{50}{Hz}, pero
\SI{9.3}{mm} a \SI{20}{Hz} (\SI{18.6}{mm} pico a pico), más que la carrera de muchos
vibradores pequeños. El anexo no fija la banda; la de \SIrange{50}{500}{Hz} es una
elección del equipo. La fuerza
pico, con CanSat, fijación y armadura (unos \SI{1.1}{kg}), es de \SI{162}{N}. Es un
tamaño habitual en los laboratorios de vibraciones de ingeniería mecánica.
\begin{enumerate}
  \item Búsqueda de resonancias: barrido senoidal de \SIrange{10}{1000}{Hz} a
        \SI{0.5}{g}, en el eje axial y en uno lateral.
  \item Exposición: barrido a 15\,G de pico entre \SIlist{50;500}{Hz}, a una octava por
        minuto, con el CanSat encendido y transmitiendo.
  \item Repetición del barrido de baja amplitud: si una frecuencia se corre más de 5\,\%,
        algo se aflojó.
\end{enumerate}
La primera frecuencia de flexión de una pala plegada, estimada como viga de placa curvada
($E\approx\SI{50}{GPa}$) apoyada en la bisagra y en la traba de la punta, es de unos
\SI{340}{Hz}: cae dentro de la banda de exposición, así que la punta y la traba verán
una respuesta amplificada y conviene un acelerómetro pequeño cerca de una punta. Si la
traba no sujeta la punta, la pala queda articulada en la bisagra y libre en el otro
extremo: no tiene un modo elástico bajo, sino que pivota y golpea contra el cuerpo. Eso
aparece en el barrido como impactos y ruido de banda ancha, y delata una traba floja.

Como alternativa sin vibrador, la guía de la competencia CanSat propone una lijadora
orbital aleatoria fijada boca arriba: de \SIrange{200}{233}{Hz} y de 20 a 29\,G,
encendida y apagada 5 veces~\cite{ensayos-cansatenv}. En ese caso los G deben
confirmarse con un acelerómetro de \SI{50}{g} o más en la fijación.

\subsection{Choque de 30\,G (E5b)}
La guía de la competencia usa una caída con cuerda de \SI{61}{cm} atada a un cáncamo, que
genera unos 30\,G~\cite{ensayos-cansatenv}. Con caída libre $L$ y rigidez equivalente $k$
de toda la cadena, el pico vale
\begin{equation}
  G = 1+\sqrt{1+\frac{2kL}{mg}}.
  \label{eq:ensayos-G}
\end{equation}
Para 30\,G hace falta $k\approx\SI{3.4}{kN/m}$, con un pulso de \SI{38}{ms}
(figura~\ref{fig:ensayos-G}). Una cuerda de Kevlar sola es mucho más rígida y daría cientos de G en
teoría; nudos, cáncamo y estructura bajan el pico real. Por eso \textbf{el pico
se mide}, con un acelerómetro de \SI{\pm200}{g} a \SI{1}{kHz} o más, porque la IMU de
vuelo ($\pm16$\,g) satura. Si el pico supera 40\,G, se intercala un tramo de cuerda
elástica, apartándose de la guía, que pide una cuerda no elástica; la desviación se
documenta. Se repite 3 veces en el eje axial. Criterio: nada se suelta, la traba sigue
cerrada (\Req{M3}), no hay reinicio ni corte de alimentación (\Req{M4}) y la telemetría
sigue.

\subsection{Térmico, radio y autonomía (E7--E9)}
\begin{description}[leftmargin=0pt, itemsep=3pt]
  \item[E7 Térmico.] Estufa de laboratorio, o caja aislante con secador y termostato, a
        \SI{60}{\celsius} durante \SI{2}{h}~\cite{ensayos-cansatenv}. El PETG ablanda cerca de
        \SI{80}{\celsius}, así que se vigilan cubo, anillo de traba y portapalas.
        \emph{Criterio}: funciona durante y después; el servo libera; el rotor gira libre más
        de \SI{5}{s}; el anillo se deforma menos de \SI{0.2}{mm}.
  \item[E8 Alcance de radio.] A \SI{915}{MHz} la pérdida en espacio libre es de
        \SI{91.7}{dB} a \SI{1}{km} y \SI{97.7}{dB} a \SI{2}{km}. Se ensaya a \SI{2}{km} con
        línea de vista y antenas elevadas (la primera zona de Fresnel tiene \SI{12.8}{m} de
        radio a mitad de camino), o a \SI{200}{m} con un atenuador de \SI{20}{dB}, que
        equivale a 10 veces la distancia. \emph{Criterio}: menos de 1\,\% de paquetes
        perdidos en \SI{10}{min} a \SI{1}{Hz}, con el CanSat en configuración de vuelo.
  \item[E9 Autonomía en plataforma.] \SI{3}{h} en \texttt{LAUNCH\_PAD} con telemetría
        (\Req{E7} pide 2) y luego la secuencia completa en modo simulación, con actuación del
        servo; se registran tensión y corriente (\Req{SN3}, \Req{SN8}). La estimación es de
        \SI{3.7}{h} con todo encendido y \SI{13}{h} con las cámaras apagadas en plataforma
        (tabla~\ref{tab:energia}). \emph{Criterio}: completa la secuencia con la celda
        sobre \SI{3.5}{V}.
\end{description}

%=====================================================================
\section{Matriz de trazabilidad}\label{sec:ensayos-matriz}
Cada requisito del anexo~\cite{anexo} tiene un método y un criterio. Métodos:
\textbf{I} inspección, \textbf{A} análisis, \textbf{E} ensayo, \textbf{D} demostración.
Ensayos E1--E6 del capítulo~\ref{cap:riesgos}; E2b, E7--E9 de este capítulo;
\textbf{E10} inspección física (masa, medidas, ajuste en un tubo de referencia);
\textbf{E11} pruebas funcionales de software y estación terrena. Los ensayos se escriben en
redonda (E7) y los requisitos del anexo en sans serif (\Req{E7}).

{\footnotesize\renewcommand{\arraystretch}{1.05}
\begin{longtable}{@{}>{\raggedright\arraybackslash}p{1.95cm}p{0.9cm}p{1.75cm}p{10.5cm}@{}}
\caption{Matriz de trazabilidad requisito $\rightarrow$ método $\rightarrow$ ensayo
$\rightarrow$ criterio.}\label{tab:ensayos-matriz}\\
\toprule
\textbf{Req.} & \textbf{Mét.} & \textbf{Ensayo} & \textbf{Criterio de aceptación}\\\midrule
\endfirsthead
\toprule
\textbf{Req.} & \textbf{Mét.} & \textbf{Ensayo} & \textbf{Criterio de aceptación}\\\midrule
\endhead
\bottomrule
\endfoot
\Req{C1} & I & E10 & Contenedor entra en una maqueta del fuselaje con las medidas de la organización\\
\Req{C2} & I, D & E10 & La expulsión la da el cohete; el CanSat no necesita orden propia y sale libre del tubo de referencia (ver \Req{S7})\\
\Req{C3} & A, E & E4 & Drogue inflado en menos de \SI{1}{s} sin orden del software; velocidad de la etapa 1 entre \SIlist{12;18}{m/s} (previsto 13,4) en caídas con rotor trabado\\
\Req{C4} & A, E & E3, E4 & $\CR\ge1{,}0$ (cota inferior 95\,\%); $\Vr$ con cota superior 95\,\% $\le\SI{7.5}{m/s}$ (margen del equipo; el anexo admite 8); media de la etapa 2 por análisis\\
\Req{C5}, \Req{X5} & D & E6, E8 & Paquetes cada \SI{1.0 \pm 0.1}{s} durante toda la secuencia\\
\Req{C6}, \Req{C7} & D & E4, E6 & Video cenit y nadir continuo desde la suelta hasta el suelo; se ven drogue, apertura y giro de palas\\
\Req{C8}, \Req{E8}, \Req{E9} & I, D & E10 & Baliza suena con la batería principal desconectada; interruptor propio accesible; audible a \SI{30}{m} (criterio del equipo)\\
\Req{C9} & I & --- & Lista de costos con comprobantes $<\$1.000.000$\\
\Req{S1} & I & E10 & Balanza de \SI{0.1}{g}: \SI{500 \pm 10}{g} (objetivo del equipo $\pm\SI{5}{g}$)\\
\Req{S2} & I & E10 & Diámetro de \SI{95}{mm} medido en 3 secciones y 2 orientaciones; las palas plegadas no sobresalen\\
\Req{S3} & I & E10 & Largo total entre \SIlist{200;250}{mm} con drogue y rotor plegados\\
\Req{S4} & E & E5a & 15\,G de pico, barrido \SIrange{50}{500}{Hz}: sin daños; resonancias corridas menos de 5\,\%\\
\Req{S5} & E & E5b & Pico medido de 30\,G o más (3 veces): sin daños, traba cerrada, sin reinicio\\
\Req{S6} & I, A, E & E10, E5 & Pared de \SI{2}{mm} o más medida en 6 puntos; sin deformación tras E5 y E7\\
\Req{S7} & D & E10 & Sale por su propio peso de un tubo vertical de referencia, 10 de 10 veces\\
\Req{S8} & I, E & E10, E5 & Todo fijado con tornillos, separadores o adhesivo; nada se mueve tras E5\\
\Req{S9} & I & --- & Lista de materiales sin contaminantes\\
\Req{M1}, \Req{M2} & I & --- & Liberación con servo; sin pirotecnia ni elementos calientes\\
\Req{M3} & E & E1, E5 & 20 de 20 liberaciones; la traba no se abre con vibración ni choque\\
\Req{M4} & I, E & E5b & Batería con conexiones soldadas o atornilladas; sin caída de tensión ni reinicio en el choque\\
\Req{E1}, \Req{E2} & I & E10 & Celda 18650 de litio con envase metálico\\
\Req{E3}, \Req{X7} & D & E11 & El panel iluminado cambia la señal transmitida; con telecomando se recibe el estado del panel y de la batería\\
\Req{E4}--\Req{E6} & I, D & E10 & Interruptores accesibles con el contenedor puesto; orificios de \SI{10}{mm} o menos; indicador visible\\
\Req{E7} & E & E9 & \SI{3}{h} en plataforma y secuencia completa después\\
\Req{X1}--\Req{X3} & I, D & E11 & Frecuencia en banda; NetID/PanID igual al número de equipo; un tercer módulo con otra dirección no recibe tramas\\
\Req{X4} & D & E4, E6 & La telemetría arranca sola al detectar la separación\\
\Req{X6} & D & E11 & El CSV contiene todos los campos con el formato y la resolución del anexo\\
\Req{SN1} & E & E4, E11 & Altitud barométrica frente a la del drone o la de una escalera medida: $\pm\SI{1}{m}$\\
\Req{SN2}, \Req{SN3}, \Req{SN8} & E & E7, E11 & Temperatura $\pm\SI{1}{\celsius}$; tensión $\pm\SI{0.05}{V}$; corriente $\pm\SI{0.01}{A}$ frente a instrumentos de referencia\\
\Req{SN4} & D & E11 & Posición GPS y número de satélites coherentes con un punto conocido\\
\Req{SN5} & E & E2, E11 & Giróscopo frente al giro conocido de una mesa giratoria; acelerómetro en reposo $\pm0{,}05$\,g\\
\Req{SN6}, \Req{SN7} & I, D & E4 & Video en color de $640\times480$ o más; se ve el despliegue\\
\Req{G1}, \Req{F5} & D & E11 & Tras \texttt{CAL} la altitud es $0\pm\SI{0.5}{m}$; un reinicio forzado conserva calibración y estado\\
\Req{G2}, \Req{G14} & D & E11 & \texttt{Flight\_<ID>.csv} con encabezado en disco y respaldo; el conteo de paquetes recibidos coincide con el registro de a bordo\\
\Req{G3}, \Req{F2}, \Req{F3} & D & E11 & Tiempo de misión con resolución de \SI{1}{s}, a menos de \SI{1}{s} de UTC tras \texttt{ST}, y conservado ante reinicio\\
\Req{G4}--\Req{G8}, \Req{G13} & I, D & E11 & GUI propia, sin pestañas, unidades SI, gráficos y campos del anexo en tiempo real\\
\Req{G9}--\Req{G12} & I, D & E11 & \SI{2}{h} con batería sin red; antena de mano o de mesa; fuente de 14\,pt o más, fondo claro, legible al sol\\
\Req{G15}, \Req{F4} & D & E11 & Todos los comandos (\texttt{CX}, \texttt{ST}, \texttt{CAL}, \texttt{MEC}, propios) responden y están documentados\\
\Req{F1} & D & E11 & El contador de paquetes sobrevive a 3 reinicios forzados sin volver a cero\\
\end{longtable}
}

%=====================================================================
\section{Hojas de registro}\label{sec:ensayos-hojas}
Cada corrida se registra en papel además del archivo digital. El nombre del archivo
lleva el código de la corrida (ej. \texttt{E3-P3-V065-R2}). Las plantillas siguientes son
las de los dos ensayos que definen \Req{C4}. Antes de cada sesión se revisan pasadores,
portapalas y destorcedor, se balancean las palas a $\pm\SI{0.1}{g}$, se prueba la traba 3
veces, se sincroniza la hora (\texttt{ST}) y se pone a cero el sensor diferencial sin
flujo. Cada anomalía se anota y alimenta el AMFE (tabla~\ref{tab:amfe}).

\providecommand{\ensayosCampo}[1]{\rule{#1}{0.4pt}}
\begin{table}[H]
\centering\small
\caption{Plantilla de registro del banco (E2/E3), una fila por punto.}
\label{tab:ensayos-hojaE3}
\begin{tabularx}{\textwidth}{@{}|Y|Y|Y|Y|@{}}
\hline
Fecha: \ensayosCampo{1.6cm} & Operadores: \ensayosCampo{1.6cm} & $p_\text{amb}$ [hPa]: \ensayosCampo{0.9cm} & $T_\text{amb}$ [\si{\celsius}]: \ensayosCampo{0.9cm}\\\hline
Paso: \ensayosCampo{1.2cm}\,° & Palas n.º: \ensayosCampo{1.4cm} & $K$: \ensayosCampo{1.2cm} & $F_b$: \ensayosCampo{1.2cm}\\\hline
\end{tabularx}\\[2pt]
\begin{tabularx}{\textwidth}{@{}|c|c|c|c|c|c|c|c|Y|@{}}
\hline
\textbf{Corr.} & \textbf{Var. [\%]} & $\boldsymbol{q_\text{ref}}$ [Pa] & $\boldsymbol{T}$ [N] & \textbf{Tara} [N] & \textbf{rpm hall} & \textbf{rpm ópt.} & $\boldsymbol{\Delta z}$ [mm] & \textbf{Notas}\\\hline
 & & & & & & & & \\[3pt]\hline
 & & & & & & & & \\[3pt]\hline
 & & & & & & & & \\[3pt]\hline
 & & & & & & & & \\[3pt]\hline
\end{tabularx}
\end{table}

\begin{table}[H]
\centering\small
\caption{Plantilla de registro de caídas desde drone (E4), una fila por caída.}
\label{tab:ensayos-hojaE4}
\begin{tabularx}{\textwidth}{@{}|Y|Y|Y|Y|@{}}
\hline
Fecha y lugar: \ensayosCampo{1.3cm} & Operador: \ensayosCampo{1.6cm} & Masa [g]: \ensayosCampo{1.2cm} & Paso: \ensayosCampo{1.2cm}\,°\\\hline
\end{tabularx}\\[2pt]
\begin{tabularx}{\textwidth}{@{}|c|c|c|c|c|c|c|c|Y|@{}}
\hline
\textbf{N.º} & \textbf{Hora} & \textbf{Caso} & \textbf{$h_0$} [m] & \textbf{Viento} [m/s] & $\boldsymbol{T_\text{loc}}$ [\si{\celsius}] & \textbf{Ventana} [s] & $\boldsymbol{V_\text{ISA}}$ [m/s] & \textbf{Notas}\\\hline
 & & & & & & & & \\[3pt]\hline
 & & & & & & & & \\[3pt]\hline
 & & & & & & & & \\[3pt]\hline
 & & & & & & & & \\[3pt]\hline
\end{tabularx}
\end{table}
